\documentclass[11pt,letterpaper]{article}
\usepackage[T1]{fontenc}
\usepackage[utf8]{inputenc}
\usepackage[scaled=0.95]{helvet}
\renewcommand{\familydefault}{\sfdefault}
\usepackage[top=1.7cm,bottom=1.5cm,left=2.2cm,right=2.2cm]{geometry}
\usepackage{microtype}
\usepackage[hidelinks]{hyperref}
\pagestyle{empty}
\setlength{\parindent}{0pt}
\setlength{\parskip}{7pt}

\begin{document}

{\large\bfseries Grover Aruquipa}\par\vspace{-6pt}
{\small Qu\'ebec, QC, Canada \ \textbullet\ \ grover.aruquipa.9@gmail.com \ \textbullet\ \ groveraruquipa.github.io}

\vspace{4pt}
Hiring Team, Humanoid, Vancouver\\
Mechanical Design Engineer, Robotic Hand

Dear Hiring Team,

A robot hand is a set of linkages that has to fit in a very small volume, move fast, hold without crushing, and survive thousands of grasps. That is the problem I have spent my PhD on, and I would like to bring it to the hand of HMND~01.

I am finishing my PhD at Universit\'e Laval with Prof.\ Cl\'ement Gosselin, whose lab has designed robotic hands and grippers for over twenty years, and Prof.\ Patrice Lambert. I defend in early 2027. During these three years I designed, built and tested three robots in which the gripper is driven through linkages and concentric shafts from motors at the base. For each one I did the full loop myself: linkage synthesis, statics and kinematics, SolidWorks design, choice of materials, motors and bearings, 3D printing and machining, assembly, integration with the electronics and control, and testing on the bench until it worked.

Two of them are close to what your hand team builds. The first is a 4-DoF coaxial wrist-gripper that sends four actuation channels through concentric hollow shafts, with 0.25\,kg of moving mass and a non-backdrivable jaw. It turned its tool ten full revolutions with less than one degree of tracking error while holding the object. The second, which I am building now, is a Delta robot with an integrated Sarrus-linkage gripper that estimates the grasp force from the motor currents, without any sensor in the fingers, so it can pick eggs quickly without breaking them. I also wrote a published review in the ASME Journal of Mechanisms and Robotics on parallel mechanisms for grasping, including parallel and underactuated robot hands, so I know the design space well and the trade-offs that come with each choice.

What I would bring to the team is speed from idea to tested hardware. I design for the process that will make the part, I iterate on the bench rather than on slides, and I am used to working closely with electrical and control engineers because on my own robots I was all three. I also lead two graduate students day to day on an 8-actuator robot, so I know how to split a mechanism into parts others can own.

I will be eligible for an open Post-Graduation Work Permit as soon as I complete my PhD, so I do not need sponsorship, and I am ready to move to Vancouver and work on site from March 2027. Prof.\ Gosselin, Prof.\ Lambert and Univ.-Prof.\ Andreas M\"uller of JKU Linz can speak to my work; their details are in my CV.

Thank you for your time. I would be glad to show you the prototypes on video or in person.

Sincerely,\\
Grover Aruquipa

\end{document}
